\scanpage{23}\mspage{22}
Il nous fallait plutôt $\mathrm b\Rightarrow\mathrm b'$ pour achever la dém. de la prop., \deleted{\unsure{Voici} \unsure{donc} \unsure{le} déjà \unsure{fait}} \inserted{et en fait on l'a déjà} \deleted{donc \unsure{à|:} voir \unsure{(pour}} \deleted{(pour $\mathrm b\Rightarrow\mathrm b'$) que}
\deleted{moyennant b), \unsure{on}}
\[ \deleted{A\in F,\quad B\in G,\quad A^\circ\subset B\quad\Longrightarrow\quad A\subset B} \]
\deleted{\unsure{On}} prouvé \deleted{\unsure{cel|cela}}en démontrant $\mathrm b)\Rightarrow\mathrm a)$ (Cf lemme marginal p. 20) Reste : prouver $\mathrm b'\Rightarrow\mathrm b$ donc moyennant $\mathrm{b}'$, le fait que $\varphi$ est croissant
\[ A,A'\in F,\quad B,B'\in G,\quad A^\circ\subset B^\circ,\quad A'^\circ\unsure{\subset|\subset'}B'^\circ,\quad \deleted{\Longrightarrow}\inserted{A'\subset A\deleted{{}^{\unsure{\circ|'|*}}}}\quad\Longrightarrow\quad B'\subset B\deleted{{}^{\unsure{\circ|'|*}}} \]
Mais \deleted{\unsure{Comme}} \inserted{par $\mathrm{b}'$) on a} \deleted{$\unsure{A^\circ\subset A^\circ|A^\circ\subset B^\circ}$ donc} \deleted{$A\subset B$, $A'\subset B'$} \deleted{$A'^\circ\subset B'^\circ$, donc par} $A\deleted{{}^{\unsure{\circ}}}\subset B\deleted{{}^{\unsure{\circ}}}$ \deleted{$\mathrm{b}'$) $A'^\circ\subset B'$, \unsure{donc} on a $A'^\circ\subset A'\subset A\subset B$ donc $A'^\circ\subset$} \deleted{$A^\circ\subset B\cap B'$, donc} donc $A'^\circ\subset B'^\circ$ \deleted{$\unsure{\subset B^\circ}\ \unsure{A\subset A'}$}
$A'^\circ\subset A\subset B$, donc $A'^\circ\subset B'^\circ\cap B$, donc $B'^\circ\cap B\ne\varnothing$, donc $B'\deleted{{}^{\unsure{\circ|'|*}}}\unsure{\in|\subset}B$, qed.

\uline{Définition} On dit, sous ces conditions, que $F$ est un \uline{raffinement} de $G$.
% slanted
\margin{\uline{NB} $F$ raffine $G$ \unsure{ssi|sss} $\forall A\in F$, $F_A$ \unsure{raff.|raffine} $G$ i.e. tt multistrate de $F$ raffine $G$.}

\uline{Ex} Toute sous figure de $G$ est un \unsure{raff.}$^{\unsure{t}}$ de $G$
Notation $F\ll G$. C'est \uline{une relation d'ordre}

\uline{Proposition} Soient $F$ et $G$ deux figures, \deleted{telles que $F$ soit un raffinement de $G$.} Conditions équivalentes
a) \inserted{$F\ll G$,} Il n'existe pas de raffinement de $G$ autre que $F$, parmi ceux dont $F$ soit une sous figure -- i.e. $F$ est \unsure{minimal|maximal} dans l'ens des raffinements de $G$, ordonné \unsure{par|pour} $\leq$ (relation de sous figure)
b) \inserted{$F\ll G$ et} $|F|=|G|$ \deleted{c) L'applicat}
